\begin{table}[H]
  \centering
  \caption{Número de cabos e massa do moitão}
  \begin{tabularx}{\textwidth}{>{\centering\arraybackslash}X>{\centering\arraybackslash}X>{\centering\arraybackslash}X}
    \toprule
    Parâmetro       & Símbolo             & Valor               \\
    \midrule
    Número de cabos & $n_c$               & $4$                 \\
    Massa do moitão & $M_{\text{moitão}}$ & \SI{75}{\kilo\gram} \\
    \bottomrule
  \end{tabularx}
\end{table}

O peso do moitão é dado pela equação:
\begin{equation}
  G_{\text{moitão}} = M_{\text{moitão}} \cdot g
  \label{eq:peso_moitao}
\end{equation}
Onde:
\begin{itemize}[leftmargin=1.25cm]
  \item $G_{\text{moitão}}$: peso do moitão em $\si{\newton}$.
  \item $M_{\text{moitão}}$: massa do moitão em $\si{\kilo\gram}$.
  \item $g$: aceleração da gravidade em $\si{\meter/\second^2}$.
\end{itemize}
\vspace{1em}

Substituindo os valores na equação \eqref{eq:peso_moitao}:
\begin{align*}
  G_{\text{moitão}} & = 75 \cdot 10               \\
  G_{\text{moitão}} & = \boxed{\SI{750}{\newton}}
\end{align*}
